\section{Experimental Validation and Testing}
\label{sec:validation}

Section~\ref{sec:system-design} explained how PyTrade is built and why each design choice was made; this chapter reports what happens when the resulting agents are actually run on data they have never been trained on. Every comparison below follows directly from a claim made in Section~\ref{sec:system-design}. The excess-return reward design (Section~\ref{sec:excess-return}) is only meaningful if the agent it trains is actually shown to beat the passive counterfactual it is scored against, not merely a strawman that is easy to beat. The shadow-portfolio mechanism that injects the Single-Asset Agent's (SAA) signal into the Portfolio Allocator Agent (PAA, Section~\ref{sec:info-flow}) is only justified if removing or corrupting that signal measurably hurts the PAA. The modular decomposition itself (Section~\ref{sec:overview}) is only a defensible answer to this thesis's research question if the resulting two-layer system outperforms both a signal-free cross-sectional model and, budget permitting, a monolithic single-network alternative trained end to end. Baselines throughout this chapter are therefore chosen to isolate one specific claim each, never to maximize the number of comparison lines on a chart, and the purged, block-level out-of-sample protocol recalled in Section~\ref{sec:oos-protocol} is what makes any of these comparisons trustworthy rather than an artifact of in-sample memorization.

Two further terms recur throughout this chapter and are fixed here once. An \emph{equity curve} is simply a portfolio's marked-to-market value plotted against time; every comparison in this chapter is, at bottom, a comparison of equity curves against some passive alternative. An \emph{ablation} is a controlled comparison that removes or replaces exactly one component of a system while holding everything else, architecture, training budget, and evaluation protocol, fixed, so that any resulting change in performance can be attributed to that one component alone. Sections~\ref{sec:paa-comparisons} and~\ref{sec:paa-implementation} rely on this idea repeatedly: the SAA's signal is the one component being switched on, off, or replaced with noise across three otherwise identical PAA variants.

This chapter is organized in the order its two agents were introduced. Section~\ref{sec:oos-protocol} recaps the evaluation protocol shared by both agents and introduces the validation periods used. Section~\ref{sec:saa-results} reports SAA results on the validation periods. Section~\ref{sec:paa-results} reports the PAA results on both validation and test periods.

\subsection{Out-of-Sample Protocol}
\label{sec:oos-protocol}

Every result in this chapter is reported on the purged, block-level split introduced in Section~\ref{sec:simulator}: repeated purged train/validation cycles carved out of the full 2000--2025 history, each separated from its own preceding training block by a 60-trading-day purge gap, followed by one held-out test block preceded by its own purge gap. Applied to the current dataset, this scheme yields four validation blocks, each spanning a distinct historical regime, and one test block reserved for a single, final evaluation. The distinction between these two evaluation surfaces matters for how the rest of this chapter should be read. The four validation blocks are used iteratively, both for the "peak alpha" checkpoint selection described in Section~\ref{sec:saa-robustness} and for every diagnostic plot in this chapter; because they are consulted repeatedly during development, a result reported only on a validation block is evidence of learning, not yet evidence of genuine, unseen-data skill. The test block is touched exactly once, after every design and checkpoint decision has already been made on the validation blocks alone, which is precisely what separates a legitimate out-of-sample result from an inflated, cherry-picked headline number.


Table~\ref{tab:oos-blocks} names each validation block by its calendar date range and summarizes the macroeconomic and financial-market conditions realized within it. This step is necessary because a purged split only guarantees that training and evaluation windows do not overlap in time; it does not guarantee that the resulting blocks sample comparable market regimes. An agent evaluated only on calm, trending markets can appear skillful without having learned genuine risk management, since correctly timing entries and exits is easier when volatility and drawdowns are historically low. Conversely, an agent evaluated on a block dominated by a single macroeconomic shock can be penalized for defensive behavior, such as temporarily reducing exposure, that is rational given that shock but is scored harshly by a reward function tuned to reward decisiveness. Naming the realized regime of every block turns this risk into an explicit, auditable property of the evaluation protocol rather than a hidden confound.

Four technical terms recur in Table~\ref{tab:oos-blocks} and are defined here once. \emph{Quantitative easing} (QE) is a central bank's large-scale purchase of government or other securities, undertaken once its short-term policy rate is already near zero, with the aim of lowering long-term interest rates; it matters for an evaluation window because it can suppress volatility and support asset prices independent of the underlying economic performance of any single firm or country. A \emph{yield-curve inversion} occurs when short-maturity government bonds yield more than long-maturity bonds, the reverse of the usual relationship; it is a widely cited leading indicator of recession, because it reflects an expectation among bond investors that the central bank will need to cut rates in the future in response to a weakening economy. The \emph{VIX} is an index of the market's expectation of near-term S\&P 500 volatility, implied by option prices; a sharp rise in the VIX signals a sudden increase in perceived risk and is commonly referred to as a volatility spike. A \emph{carry trade} is a strategy of borrowing in a currency with a low interest rate to invest in higher-yielding assets denominated in another currency; an \emph{unwind} of such a trade, triggered for instance by a rise in the funding currency's interest rate, forces borrowers to close these positions quickly, producing fast, correlated selling across otherwise unrelated asset classes.

The test block (final row of Table~\ref{tab:oos-blocks}) is deliberately not characterized narratively in advance of evaluation. Naming its dominant events before the single, final evaluation pass described above would risk tempting a post-hoc, narrative-driven justification of whatever result that pass produces, which is precisely what the purged protocol is designed to prevent. Its realized composition is instead reported alongside its results.

\begin{table}[p]
\centering
\small
\begin{tabular}{@{}p{2.46cm}p{6.85cm}p{6.85cm}@{}}
\toprule
\textbf{Block (dates)} & \textbf{Realized economic and financial-market conditions} & \textbf{Relevance to agent training and evaluation} \\
\midrule
val\_00 \newline (2005-12-16 to \newline 2007-02-14) &
The U.S. policy rate stood at a cyclical peak of 5.25\%. The U.S. yield curve was inverted for much of 2006. U.S. and European equities advanced together at low realized volatility, a period later called the "Great Moderation". Asian equities advanced alongside strong global trade. Oil and gold both trended upward. Late in the block, U.S. subprime-mortgage delinquencies began rising, and a major subprime lender reported large losses. Prices had not yet reacted by the block's end. &
Nearly every instrument trends in the same direction at low volatility. This is the regime least likely to penalize a non-diversifying policy. The block ends just before, not during, a realized shock. Comparing performance here against val\_02 and val\_03 helps separate general risk management from a fit to one calm regime (Section~\ref{sec:saa-robustness}). \\
\addlinespace
val\_01 \newline (2012-01-12 to \newline 2013-03-20) &
European sovereign-debt stress continued. Greek government debt was restructured in March 2012, and Spanish and Italian bond yields stayed elevated. European equities moved with this stress. The European Central Bank pledged in July 2012 to preserve the euro. The Federal Reserve announced a third round of asset purchases (QE3) in September 2012. Gold traded near record levels. Asian equities recorded inflows attributed to a search for yield. Oil stayed elevated on supply concerns linked to sanctions on Iran. &
Policy announcements, not firm fundamentals, dominate this block. European equities move together on crisis news, largely apart from Asian equities and commodities. A checkpoint whose relative ranking across instruments differs sharply from val\_00 is responding to regime, not repeating one fixed allocation. \\
\addlinespace
val\_02 \newline (2018-02-27 to \newline 2019-04-23) &
The block opens just after a sharp volatility spike and the collapse of several short-volatility products in February 2018. The U.S. and China then imposed successive tariffs on each other. The Federal Reserve raised its policy rate four times. Export-oriented Asian and European equities moved with the tariff dispute; UK equities also traded through separate withdrawal negotiations from the European Union. Equities fell sharply together in December 2018. Oil fell over the same weeks, while gold rose. From January 2019, after the Fed signaled a pause, equities recovered. &
This block contains a decline and a recovery inside one window. It is where equity-oil co-movement is clearest, with gold moving the other way. Tests whether the agent can shift from a defensive posture back to a risk-seeking one as conditions improve, rather than only having learned to reduce exposure. \\
\addlinespace
val\_03 \newline (2024-02-26 to \newline 2025-04-22) &
U.S. equity gains concentrated in a few large technology firms. Taiwanese equities, weighted toward semiconductor manufacturing, moved closely with this trend. The Federal Reserve began cutting its policy rate in September 2024. In August 2024, following a Bank of Japan rate increase, Japanese equities recorded their largest single-day decline since 1987, attributed to an unwind of yen-funded carry positions; other equities fell by a smaller amount the same day. Gold reached repeated record levels. In April 2025, new U.S. tariffs were followed by a broad, short-lived decline across nearly all equities and in oil. &
This block contains two shocks of different origin: a single-country, rate-driven decline, and a broad, trade-policy-driven decline. It also contains a narrow, concentration-driven rally that a diversified allocation is not automatically exposed to. Comparing the response to each shock tests whether risk behavior generalizes, rather than having been fit to one event (Section~\ref{sec:saa-robustness}). \\
\addlinespace
test\_00 \newline (2025-07-21 to \newline 2026-08-06) &
Not characterized in advance of evaluation, for the reason given above. &
Reserved for the single, final evaluation described later in this section. \\
\bottomrule
\end{tabular}
\caption{Calendar date ranges of the four validation blocks and the single test block, the market conditions observable within each, and their role in selecting a checkpoint (Section~\ref{sec:saa-robustness}).}
\label{tab:oos-blocks}
\end{table}


\subsection{Single-Asset Agent Results}
\label{sec:saa-results}

This section reports results for the SAA across all four validation blocks. Three questions structure the presentation: does the agent do anything other than imitate a passive strategy, does it do so consistently across assets and time, and is whatever skill it shows attributable to genuine learning rather than to memorizing the training blocks? Every SAA chart in this chapter plots three lines against the same time axis: a \emph{cash-only} baseline, which never enters a position and is exposed only to the risk-free carry described in Section~\ref{sec:saa}; a \emph{buy-and-hold} baseline, which purchases the same asset the agent is trading at the start of the episode and never trades again, net of the same transaction-cost model described in Section~\ref{sec:friction}; and the trained SAA agent itself. The aggregated figures also show SPY performance as direct comparison. 

\subsubsection{Per-Asset, Per-Validation-Episode Detail}
\label{sec:saa-per-asset}

Figure~\ref{fig:saa-per-asset} shows one equity-curve panel per asset for each of the four validation episodes, already computed at the time of writing. Each panel plots the three baselines of Section~\ref{sec:saa-baselines} over the full length of that validation block. Two questions are asked of every panel. First, does the agent track or exceed buy-and-hold specifically during the sub-periods within the episode that are trending versus choppy, since a purely reactive agent is expected to struggle exactly where a trend reverses quickly. Second, does the agent visibly avoid the worst of buy-and-hold's drawdowns, which would be the clearest sign that its excess-return and drawdown-penalty reward terms (Sections~\ref{sec:excess-return} and~\ref{sec:drawdown}) are shaping genuinely risk-aware behavior rather than merely a noisier version of buy-and-hold.


\subsubsection{Aggregated Cross-Asset View versus SPY}
\label{sec:saa-aggregate}

\subsubsection{Training Diagnostics from TensorBoard}
\label{sec:saa-diagnostics}


\subsubsection{Statistical Robustness}
\label{sec:saa-robustness}

Checkpoint selection is a search, and every search over a finite set of validation blocks risks selecting a checkpoint that merely happened to fit those specific blocks well. This section reports the "peak alpha" checkpoint, defined here as the training checkpoint whose combined, cross-asset outperformance over SPY (Section~\ref{sec:saa-aggregate}), averaged across the four validation blocks, is highest across all checkpoints saved during training. Reporting a single peak checkpoint is only informative alongside the shape of the curve around it, so this section also reports the \emph{validation-decay} curve: the same aggregated validation metric plotted against checkpoint index, before and after the peak. A validation-decay curve that falls away sharply and monotonically after the peak is evidence that later checkpoints are overfitting to the training blocks, exactly the pattern that Section~\ref{sec:saa-diagnostics}'s training-versus-validation reward divergence would independently predict. Reading these two diagnostics together, rather than either in isolation, is what turns "the agent beat SPY" into a defensible, rather than a merely fortunate, empirical claim.

\subsection{Portfolio Allocator Agent Results}
\label{sec:paa-results}


\subsubsection{Implementation Plan for the Three New Agents}
\label{sec:paa-implementation}

The \textbf{control ablation} never loads or runs the frozen SAA at all, at training time or at inference time, which makes it computationally cheaper than the hierarchical system rather than merely architecturally identical to it. In place of the SAA's own forward pass, the signal injected into each asset's token is drawn from an autoregressive, order-one (AR(1)) noise process, clipped to the same range as the SAA's own action output. An AR(1) process is chosen deliberately over independent, identically distributed (i.i.d.) noise: because it is autocorrelated from one day to the next, it is at least superficially smooth in the way a genuine trading signal is smooth, which is a harder, and therefore more informative, control condition than noise a network could trivially learn to ignore as pure static. Crucially, this noise is still committed to each asset's real, isolated shadow sub-portfolio through the same shadow-portfolio bookkeeping already described in Section~\ref{sec:info-flow}, so the shadow holding percentage and the other shadow-derived diagnostics injected alongside the signal remain populated exactly as they are for the hierarchical system; only the origin of the committed action changes, from a frozen policy to a stochastic process.

The \textbf{cross-sectional-only ablation} is implemented as the zero-valued special case of the same mechanism: the injected signal is held at a constant zero rather than drawn from the AR(1) process, and the corresponding shadow sub-portfolio is simply never traded. This is a deliberate design choice rather than an implementation shortcut. Implementing this ablation as a genuinely smaller network, one whose asset token omits the SAA-signal dimension entirely, would confound two different explanations for any performance gap: a loss of information, and a loss of model capacity. Holding the signal at zero instead keeps this ablation, the control ablation, and the hierarchical system at identical architecture and identical parameter count, differing only in what occupies two input columns of the asset token. This isolates the information content of the SAA's signal specifically, which is the comparison Section~\ref{sec:paa-comparisons} actually needs.

The \textbf{monolithic end-to-end baseline} is the only one of the three that must be built substantially from new code, since neither the shadow-signal injection mechanism nor the split asset/portfolio token design of Section~\ref{sec:paa-state} presuppose a frozen, per-asset temporal specialist to consume. The simulator itself, its reward, its friction model, and its purged evaluation protocol are reused entirely unchanged, as is the surrounding validation and checkpointing infrastructure, since none of that infrastructure depends on what kind of network is being trained. What is genuinely new is a single policy that supplies its own temporal memory over raw, per-asset market features and performs both timing and cross-sectional allocation jointly, without ever being told which sub-decision is which. Wherever possible, this baseline is built to reuse the hierarchical system's own architectural building blocks, for instance replacing the frozen SAA's contribution with a trainable, per-asset recurrent memory feeding directly into the same shared-attention-plus-private-heads structure described in Section~\ref{sec:paa-network}, rather than adopting an unrelated architecture from the literature. The reasoning behind this choice is methodological: if the monolithic baseline is built from wholly different architectural primitives, any performance gap between it and the hierarchical system could be attributed to that unrelated difference rather than to the modular-versus-monolithic distinction this thesis is actually testing. Training compute, measured either in total environment steps or in wall-clock time, is matched explicitly against the budget spent training the hierarchical system, so that the comparison in Section~\ref{sec:paa-comparisons} is not confounded by one system simply being trained longer than the other. Because this baseline is the most implementation-intensive of the three, it is explicitly time-boxed: if training time does not permit it before submission, its omission is reported as a limitation of this thesis (Section~\ref{sec:limitations}) rather than silently left out of this chapter.

\subsubsection{Multi-Agent Inference and Comparison Harness}
\label{sec:paa-harness}


\subsubsection{Baselines and Ablations to Compare}
\label{sec:paa-comparisons}

Table~\ref{tab:paa-comparisons} lists the full comparison set the harness of Section~\ref{sec:paa-harness} must produce, ordered from the simplest passive reference to the most demanding architectural alternative. Each row isolates one specific claim, matching this chapter's opening principle that baselines are chosen for what they isolate.

\begin{table}[htbp]
\centering
\begin{tabular}{@{}p{5cm}p{6.0cm}p{5cm}@{}}
\toprule
\textbf{Condition} & \textbf{What it isolates} & \textbf{Status} \\
\midrule
100\% cash & Floor: reward from risk-free carry alone & Available (env.\ default) \\
SPY buy-and-hold & The PAA's own reward counterfactual (Section~\ref{sec:excess-return}) & Available (env.\ default) \\
Fixed-weight benchmark & A practitioner-style static allocation, no learning at all & Available (env.\ default) \\
Cross-sectional-only ablation & Whether the SAA's temporal signal adds value at all & Not yet trained \\
Control ablation (randomized signal) & Whether any gain over the row above is genuine information, not just an extra input column & Not yet trained \\
Full hierarchical PAA & The proposed system & Trained; not yet evaluated against the rows above \\
Monolithic end-to-end baseline & The modular-versus-monolithic research question directly & Time-boxed stretch goal (Section~\ref{sec:paa-implementation}) \\
\bottomrule
\end{tabular}
\caption{The full PAA comparison set. Rows are ordered from the simplest passive reference to the most architecturally demanding alternative; each isolates exactly one claim made in Section~\ref{sec:system-design}.}
\label{tab:paa-comparisons}
\end{table}

\subsubsection{Statistical Robustness}
\label{sec:paa-robustness}

Once the comparisons of Table~\ref{tab:paa-comparisons} exist, they are subjected to the same treatment already applied to the SAA in Section~\ref{sec:saa-robustness}: a peak-alpha checkpoint is identified on the four validation blocks, a validation-decay curve is reported around that peak as evidence for or against overfitting, and both are cross-referenced against PAA-specific training diagnostics analogous to those described in Section~\ref{sec:saa-diagnostics}. Two diagnostics are specific to the PAA's own architecture and have no SAA equivalent: the stability of the attention weights the shared \texttt{TransformerEncoder} (Section~\ref{sec:paa-network}) assigns to each asset token across checkpoints, and the degree of divergence between the actor and critic losses described in Section~\ref{sec:actor-critic-theory}, which the private per-head encoder layers of Section~\ref{sec:paa-network} are specifically designed to absorb.
